\documentclass[11pt,a4paper]{article}
\usepackage[T1]{fontenc}
\usepackage[utf8]{inputenc}
\usepackage{lmodern,microtype,amsmath,amssymb,amsthm,booktabs}
\usepackage[margin=27mm]{geometry}
\usepackage[dvipsnames]{xcolor}
\usepackage[colorlinks=true,linkcolor=MidnightBlue,citecolor=MidnightBlue,urlcolor=MidnightBlue]{hyperref}
\newtheorem{theorem}{Theorem}
\newtheorem{corollary}[theorem]{Corollary}
\newtheorem{proposition}[theorem]{Proposition}
\newcommand{\E}{\mathbb E}
\newcommand{\Law}{\operatorname{Law}}
\newcommand{\Var}{\operatorname{Var}}
\newcommand{\MEC}{\operatorname{MEC}}
\newcommand{\causal}{\mathrm{causal}}
\newcommand{\ii}{\imath}
\newcommand{\N}{\mathcal N}
\newcommand{\cst}{c_0}
\title{The Square-Root Barrier for Causal Entropy\\[3pt]
\large A converse for the crash-test pair, and a constructive route beyond it}
\author{Technical addendum to \emph{Who Pays the Bill?}, Version 4}
\date{26 September 2026}
\begin{document}
\maketitle
\begin{abstract}
Two constant action policies already impose a finite-horizon entropy penalty.
For two finite distributions of equal entropy $h$, we prove
\[
\liminf_{n\to\infty}\frac{C_n^{\causal}-nh}{\sqrt n}
\ge \frac{|\sqrt{V_P}-\sqrt{V_Q}|}{\sqrt{2\pi}},
\]
where $V_P,V_Q$ are their varentropies. For the crash-test pair in Version 4,
the coefficient is approximately $0.0173736270$; logarithmic and bounded
residuals are therefore impossible for that pair. We also sharpen the
anti-diagonal construction to $O(\sqrt n)$ whenever entropy and varentropy
match, and give an explicit nonpermutation dyadic pair satisfying these
conditions. No logarithmic upper bound is claimed.
\end{abstract}

\section{The precise causal problem}
All logarithms are base two and entropies are in bits. Let $P,Q$ be finite
probability vectors; zero-probability symbols can be discarded. At each time,
an action selects the output law $P$ or $Q$. A discrete seed $S$, drawn before
the interaction, determines all responses as functions of the action and
the preceding history. For every adaptive policy, the resulting conditional
output laws must agree exactly with this memoryless channel.

Let $C_n^{\causal}(P,Q)$ be the infimum of $H(S)$ over such horizon-$n$
realizations. This is the history-indexed contract C3-C of \cite{mikulik};
no equality of responses across distinct histories is imposed.
When $H(P)=H(Q)=h$, write
\[
E_n(P,Q)=C_n^{\causal}(P,Q)-nh.
\]
For a finite law $A$, set
\[
\ii_A(x)=-\log_2 A(x),\qquad
\mu_A=\Law(\ii_A(X)),\quad X\sim A,
\]
and $V_A=\Var(\ii_A(X))$, $\sigma_A=\sqrt{V_A}$.
For probability laws on the real line, $W_1$ denotes the infimum of
$\E|U-V|$ over their couplings. Notice that
$\mu_{P^{\otimes n}}$ is the law of a sum of $n$ independent copies of
$\ii_P(X)$.

The point of the argument is that a causal seed must support both complete
constant-action paths. Those paths may be correlated, but neither can be
deleted from the counterfactual tree.

\section{A converse from two paths}
\begin{theorem}[Product information-spectrum converse]\label{thm:endpoint}
For every finite $P,Q$ and every positive integer $n$,
\begin{equation}\label{eq:endpoint}
C_n^{\causal}(P,Q)\ge
\frac{nH(P)+nH(Q)+W_1(\mu_{P^{\otimes n}},\mu_{Q^{\otimes n}})}{2}.
\end{equation}
In particular, when $H(P)=H(Q)=h$,
\begin{equation}\label{eq:residual}
E_n(P,Q)\ge \tfrac12 W_1(\mu_{P^{\otimes n}},\mu_{Q^{\otimes n}}).
\end{equation}
\end{theorem}
\begin{proof}
Fix an admissible seed. Running the policy that always chooses $P$ gives a
deterministic function $U^n(S)$ with law $P^{\otimes n}$. Running the policy
that always chooses $Q$ gives $V^n(S)$ with law $Q^{\otimes n}$.
For each seed atom $s$ of positive probability $r_s$,
\[
r_s\le P^{\otimes n}(U^n(s)),\qquad
r_s\le Q^{\otimes n}(V^n(s)).
\]
Consequently, with $A=\ii_{P^{\otimes n}}(U^n)$ and
$B=\ii_{Q^{\otimes n}}(V^n)$,
\begin{align*}
H(S)&\ge\E\max\{A,B\}\\
&=\tfrac12\bigl(nH(P)+nH(Q)+\E|A-B|\bigr)\\
&\ge\tfrac12\bigl(nH(P)+nH(Q)
  +W_1(\mu_{P^{\otimes n}},\mu_{Q^{\otimes n}})\bigr).
\end{align*}
Take the infimum over all admissible seeds. Infinite seed entropy causes
no difficulty, since the desired lower bound is then automatic.
\end{proof}

This is a block-level converse for C3-C. It does not assert the stronger
linear, single-letter lower bound available under the synchronous contract
C3-S. Information-spectrum converses for minimum-entropy couplings are an
established method; see \cite{shkel}. The proof above is included to make
the application to the causal contract explicit.

\begin{theorem}[Varentropy obstruction]\label{thm:variance}
Suppose $H(P)=H(Q)=h$. Then
\begin{equation}\label{eq:limit}
\lim_{n\to\infty}
\frac{W_1(\mu_{P^{\otimes n}},\mu_{Q^{\otimes n}})}{\sqrt n}
=\sqrt{\frac2\pi}\,|\sigma_P-\sigma_Q|.
\end{equation}
Hence
\begin{equation}\label{eq:main}
\boxed{\displaystyle
\liminf_{n\to\infty}\frac{E_n(P,Q)}{\sqrt n}
\ge \frac{|\sigma_P-\sigma_Q|}{\sqrt{2\pi}}.}
\end{equation}
In particular, $E_n=o(\sqrt n)$ requires $V_P=V_Q$.
\end{theorem}
\begin{proof}
Let $A_n=(\ii_{P^{\otimes n}}-nh)/\sqrt n$ denote the centered,
normalized random information content, and define $B_n$ similarly.
The central limit theorem gives weak convergence to
$\N(0,\sigma_P^2)$ and $\N(0,\sigma_Q^2)$. Their second moments are
the fixed constants $V_P,V_Q$, so their absolute first moments are uniformly
integrable. The convergences therefore hold in $W_1$ as well.
The triangle inequality implies convergence of $W_1(\Law(A_n),\Law(B_n))$
to the distance between these normal laws. In one dimension the quantile
coupling is optimal; coupling $\sigma_P Z$ and $\sigma_Q Z$, with
$Z\sim\N(0,1)$, gives distance
$|\sigma_P-\sigma_Q|\E|Z|=\sqrt{2/\pi}|\sigma_P-\sigma_Q|$.
Translation invariance and scaling of $W_1$ prove \eqref{eq:limit};
\eqref{eq:residual} proves \eqref{eq:main}. Degenerate normals are allowed.
\end{proof}

\paragraph{An explicit finite-$n$ version.}
If both variances are positive, put
\[
\rho_A=\E|\ii_A(X)-H(A)|^3,\qquad
B_{P,Q}=\frac12\left(\frac{\rho_P}{V_P}+\frac{\rho_Q}{V_Q}\right).
\]
Goldstein's mean central limit theorem bound \cite{goldstein} gives
\begin{equation}\label{eq:goldstein}
W_1\bigl(\Law(\ii_{A^{\otimes n}}-nH(A)),\N(0,nV_A)\bigr)
\le \frac{\rho_A}{V_A}.
\end{equation}
The reverse triangle inequality and Theorem~\ref{thm:endpoint} yield, for
every $n\ge1$,
\begin{equation}\label{eq:finite}
E_n(P,Q)\ge
\max\left\{0,\frac{|\sigma_P-\sigma_Q|}{\sqrt{2\pi}}\sqrt n-B_{P,Q}\right\}.
\end{equation}
The square-root obstruction itself does not depend on this quantitative
normal-approximation theorem.

\section{The crash-test pair: logarithmic growth is excluded}
Take the pair from \cite{mikulik},
\[
P=(1/2,1/4,1/4),\qquad Q=(q,q,1-2q),
\quad H(Q)=3/2,\quad q\in(1/3,1/2).
\]
Its entropies agree. Under $P$, the information content equals $1$ or $2$,
each with probability $1/2$, so $V_P=1/4$.
Write $r=1-2q$. Under $Q$ there are two information levels, hence
\begin{equation}\label{eq:vq}
V_Q=r(1-r)\left(\log_2\frac qr\right)^2.
\end{equation}

\paragraph{An exact certificate that $V_Q<1/4$.}
The function $H(q,q,1-2q)$ has derivative
$2\log_2((1-2q)/q)<0$ on $(1/3,1/2)$.
For $q=a/b$ with $0<2a<b$, an exact integer comparison gives
\[
H(q,q,1-2q)>\tfrac32
\quad\Longleftrightarrow\quad
b^{2b}>2^{3b}a^{4a}(b-2a)^{2(b-2a)}.
\]
The sign is positive for $(a,b)=(41,100)$ and negative for
$(a,b)=(411,1000)$. Thus $0.410<q<0.411$ and $0.178<r<0.180$.
It follows that
\[
r(1-r)<\frac{369}{2500},\qquad
\frac qr<\frac{411}{178},\qquad
411^4<32\cdot178^4.
\]
The last inequality gives $\log_2(q/r)<5/4$, and therefore
\[
\boxed{V_Q<\frac{369}{1600}<\frac14=V_P.}
\]
All the displayed integer comparisons are reproduced by the accompanying
Python script. The proof of the strict variance gap uses no rounded
logarithm or numerically located root.

\begin{corollary}[Resolution of the bounded/logarithmic alternatives]
For this crash-test pair,
\[
\liminf_{n\to\infty}\frac{E_n}{\sqrt n}\ge c,
\qquad
c=\frac{1/2-\sqrt{V_Q}}{\sqrt{2\pi}}>0.
\]
In particular $E_n\ne O(\log n)$ and $E_n\ne O(1)$.
Numerically,
\begin{align*}
q&\approx0.41003242818198557608,\\
V_Q&\approx0.20834731023907319744,\\
c&\approx0.01737362702590893204,\\
B_{P,Q}&\approx0.66879251142378039507.
\end{align*}
The exact all-$n$ lower bound is \eqref{eq:finite}, with these constants
defined by their formulas, not by the rounded decimal values.
\end{corollary}

Together with Corollary 19.2 of \cite{mikulik}, this gives the presently
established range in this addendum
\begin{equation}\label{eq:range}
\boxed{\Omega(\sqrt n)\ \le\ E_n\ \le\ O(n^{2/3}).}
\end{equation}
This excludes two alternatives in Open Problem 27.3 of that manuscript.
It does \emph{not} determine the exact asymptotic order, nor prove a
matching $O(\sqrt n)$ upper bound for the crash-test pair.

\section{A constructive improvement when varentropies match}
The same normal-approximation estimate also improves the upper bound for
a different class of pairs. The construction is explicit, although the
block alphabets can be exponentially large.

\begin{theorem}[Matched-varentropy anti-diagonal construction]\label{thm:upper}
Suppose $H(P)=H(Q)=h$ and $V_P=V_Q=V>0$. Define
\[
D=\frac{\rho_P+\rho_Q}{2V}+\cst,
\qquad \cst=\frac{\log_2 e}{e}.
\]
For $n=mL$, where $m\ge2$ and $L\ge1$ are integers, there is an exact
causal seed with
\begin{equation}\label{eq:upper}
H(S)\le nh+Lh+(m-1)D.
\end{equation}
In particular, $E_n(P,Q)=O(\sqrt n)$ for all positive integers $n$.
\end{theorem}
\begin{proof}
Since the two centered information sums have the same normal approximation,
\eqref{eq:goldstein} and the triangle inequality give
\[
W_1(\mu_{P^{\otimes L}},\mu_{Q^{\otimes L}})
\le \frac{\rho_P+\rho_Q}{V}.
\]
The two-distribution greedy/profile bound of \cite[Theorem 4.1]{compton},
together with the profile identity in \cite[Lemma 17.1]{mikulik}, gives a
block coupling $J_L\in\Pi(Q^{\otimes L},P^{\otimes L})$ satisfying
\[
H(J_L)\le Lh+
\tfrac12 W_1(\mu_{P^{\otimes L}},\mu_{Q^{\otimes L}})+\cst
\le Lh+D.
\]
For completeness, the profile identity follows by integrating
$\max\{F_{\mu_A}^{-1}(u),F_{\mu_B}^{-1}(u)\}$ over $0<u<1$ and using
$\max\{a,b\}=(a+b+|a-b|)/2$.

Apply the anti-diagonal construction of \cite[Section 18]{mikulik}.
Split two length-$n$ streams $X,W$ into $m$ blocks of length $L$.
Independently draw $X^{(1)}\sim P^{\otimes L}$,
$W^{(m)}\sim Q^{\otimes L}$, and
\[
(W^{(j)},X^{(j+1)})\sim J_L,\qquad j=1,\ldots,m-1.
\]
On the $a$th occurrence of action $P$, output $X_a$; on the $b$th
occurrence of action $Q$, output $W_{n+1-b}$.
Observing both members of any coupled pair requires at least
$jL+1$ occurrences of $P$ and $n-jL+1$ occurrences of $Q$, hence at
least $n+2$ actions. No horizon-$n$ path can do this.
For any fixed action word, exposed coordinates therefore have their
required product law. Each output branch of an adaptive deterministic
policy determines such an action word, proving exact adaptive validity;
independent policy randomization follows by conditioning.

The independent components of the seed have total entropy at most
\[
2Lh+(m-1)(Lh+D)=nh+Lh+(m-1)D.
\]
For general $n\ge3$, choose $L=\lceil\sqrt n\rceil$,
$m=\lceil n/L\rceil$, and $N=mL$. Restrict the horizon-$N$ construction
to its first $n$ actions. Since $N-n<L$ and $m-1<n/L$,
\[
E_n\le (N-n)h+Lh+(m-1)D
\le (2h+D)\sqrt n+2h.
\]
The finitely many smaller horizons do not affect the stated order.
\end{proof}

If $V_P=V_Q=0$ and the entropies agree, both distributions are uniform
on the same number of positive-probability symbols. They are permutations
of one another, and $E_n=0$ by sharing an iid uniform stream and relabeling
its outputs. Thus the zero-variance case is also completely covered.

\section{An explicit different pair for the next logarithmic attempt}
Here are actual \emph{output-symbol distributions}, rather than auxiliary
increment laws. Repeated entries represent distinct symbols:
\begin{align*}
P_\star&=\bigl(1/8,\underbrace{1/32,\ldots,1/32}_{24\ \mathrm{symbols}},
                  \underbrace{1/128,\ldots,1/128}_{16\ \mathrm{symbols}}\bigr),\\
Q_\star&=\bigl(\underbrace{1/16,\ldots,1/16}_{8\ \mathrm{symbols}},
                  \underbrace{1/64,\ldots,1/64}_{32\ \mathrm{symbols}}\bigr).
\end{align*}
Their total masses are $1$; they have $41$ and $40$ positive-probability
symbols, so they are not permutations. If a common output alphabet is
required, append one zero to $Q_\star$.
Their information laws are especially simple:
\[
\mu_{P_\star}=\tfrac18\delta_3+\tfrac34\delta_5+\tfrac18\delta_7,
\qquad
\mu_{Q_\star}=\tfrac12\delta_4+\tfrac12\delta_6.
\]
Direct rational arithmetic gives
\begin{center}
\begin{tabular}{lcc}
\toprule
Quantity & $P_\star$ & $Q_\star$\\
\midrule
Entropy & $5$ & $5$\\
Varentropy & $1$ & $1$\\
Third centered moment & $0$ & $0$\\
Fourth centered moment & $4$ & $1$\\
Third absolute centered moment & $2$ & $1$\\
\bottomrule
\end{tabular}
\end{center}
The shifted information generating polynomials satisfy the exact identity
\[
\left(\frac18+\frac34z^2+\frac18z^4\right)
-\left(\frac12z+\frac12z^3\right)=\frac{(z-1)^4}{8}.
\]
Thus the first three information moments match, while the fourth differs.
Theorem~\ref{thm:upper} gives $E_n(P_\star,Q_\star)=O(\sqrt n)$ with
$D=3/2+\cst$. For square horizons $n=L^2$, $L\ge2$, the concrete seed obeys
\[
H(S)\le5n+\left(\frac{13}{2}+\cst\right)\sqrt n
             -\left(\frac32+\cst\right).
\]
For these same marginals, the noncausal product-block problem already has
a uniform bound
\[
0\le \MEC(P_\star^{\otimes n},Q_\star^{\otimes n})-5n
\le \frac32+\cst.
\]
This last statement follows directly from the block coupling bound in the
proof. It is not a bounded-residual causal construction: a coupling of two
constant paths alone does not supply all mixed adaptive paths.

\paragraph{What remains to be constructed.}
For this new pair, a logarithmic causal residual is not excluded by the
varentropy converse. Establishing it requires a causal seed of entropy
$5n+O(\log n)$, valid under every adaptive policy. The present construction
still pays a boundary cost proportional to $L$ and a constant cost for each
of approximately $n/L$ blocks. Balancing those costs yields $\sqrt n$.
Fourth-order moment cancellation alone does not remove either cost.

\paragraph{Scope and reproducibility.}
The main lower bound is self-contained apart from the classical central
limit theorem and elementary one-dimensional transport facts. The explicit
finite-$n$ constant and the matched-varentropy upper bound use the cited
quantitative results. No historical priority claim is made for those tools
or their combination. The supplied script checks the exact integer
certificates, verifies the new pair with rational arithmetic, and computes
the displayed decimals at 80-digit working precision. It does not substitute
numerical tests for causal validity or asymptotic proofs.

\begin{thebibliography}{9}
\bibitem{mikulik}
J. Mikul\'ik.
\emph{Who Pays the Bill? Substrate Skepticism, Stable Counterfactual
Identity, Hidden Synergy, and the Causal Shannon Floor}.
Version 4, 19 August 2026. User-supplied manuscript.

\bibitem{shkel}
Y. Y. Shkel and A. K. Yadav.
\emph{Information Spectrum Converse for Minimum Entropy Couplings and
Functional Representations}. ISIT, 2023.
\url{https://arxiv.org/abs/2305.05745}.

\bibitem{goldstein}
L. Goldstein.
\emph{Bounds on the constant in the mean central limit theorem}.
Annals of Probability 38(4), 1672--1689, 2010.
\url{https://arxiv.org/abs/0906.5145}.

\bibitem{compton}
S. Compton, D. Katz, B. Qi, K. Greenewald, and M. Kocaoglu.
\emph{Minimum-Entropy Coupling Approximation Guarantees Beyond the
Majorization Barrier}. AISTATS, PMLR 206, 10445--10469, 2023.
\url{https://proceedings.mlr.press/v206/compton23a.html}.
\end{thebibliography}
\end{document}
